\section{Forecasters and pre-registered design}\label{sec:design}

\subsection{Forecasters}
All forecasters predict the 24 hourly loads of day $d$ at midnight, using loads up to the previous day and the day's temperature. F1 repeats the load of the same hour one day earlier and F2 the load of the same hour one week earlier. F3 is a ridge regression and F4 a gradient-boosted tree model (LightGBM~\cite{ke2017lightgbm}, 400 trees) on hour-of-day and day-of-week indicators, the loads 24, 48 and 168 hours earlier, the temperature, its square and heating and cooling degree terms. F5 is the same LightGBM model trained on the pinball loss of the 0.9 quantile, a peak-conservative forecaster. Models are fitted per building on 2016 and are not refitted. The temperature is the observed value, which amounts to a perfect weather forecast for all forecasters.

To separate the type of error from its size, three synthetic forecasts are built from the actual load with the same MAE as F4 in every building-month: actual load plus Gaussian noise (\emph{noise}); the actual load moved towards the load one hour earlier by a multiple of their difference, so that rises and peaks are predicted late (\emph{1-h shift}); and the actual load scaled up by a constant factor within the month (\emph{uniform over-forecast}). Matching the MAE makes that factor equal to F4's relative MAE in the month, \res{scaleMed}\% at the median (interquartile range \res{scaleLo}--\res{scaleHi}\%); the analysis plan had described this error as a 5\% scale error, but the matched-MAE construction it specified fixes the size.

Accuracy is measured per building-month by the MAE, the RMSE, the mean absolute percentage error, the absolute error at the hour of the actual daily peak, the error of the forecast daily peak magnitude and the timing error of the forecast daily peak. Over the test year the median MAE was \res{maeF4} kWh per hour for F4, \res{maeF3} for F3, \res{maeF1} for F1, \res{maeF2} for F2 and \res{maeF5} for F5 (relative to the mean load: \res{relmaeF4}\%, \res{relmaeF3}\%, \res{relmaeF1}\%, \res{relmaeF2}\% and \res{relmaeF5}\%). Within a building-month the two summaries rank the forecasters identically; ranked by MAE across building-months, F4 was most accurate on average (mean rank \res{rankF4}, best in \res{bestF4}\%) and F5 least (mean rank \res{rankF5}, last in \res{worstF5}\%). The median relative bias was close to zero for F1--F4 and \res{relbiasF5}\% for F5.

\subsection{Pre-registration and amendments}
The questions, hypotheses, primary test, forecasters, tariffs, battery and robustness checks were fixed in an analysis plan committed to the study repository (\url{\repoURL}, file \texttt{PLAN.md}) before any data were downloaded. Two amendments followed, both before the full run and both in the repository. The first, written with the code and before any data were loaded, fixed the exclusion rule for building-months with a negligible perfect-information saving and the reset of the state of charge. The second followed a pipeline check on six buildings whose outcomes were seen: the original rule of excluding any building with a gap longer than three hours kept only \res{nKeptThreeH} of the \res{nMeters} meters (\res{nExclThreeH} excluded for that reason alone), so longer gaps were filled as described above; and because every forecaster executed open loop raised the bill in those six buildings, the guarded execution was added as a secondary controller. The hypotheses and their tests were not changed, and the primary results below are those of the registered open-loop controller.

\subsection{Hypotheses and tests}
\begin{itemize}
\item[H1] (primary) Under T1, the ranking of F1--F5 by MAE agrees with their ranking by captured share. For each building-month we compute Kendall's $\tau$ between the negative MAE and the share across the five forecasters, average it over the months of each building, and test the mean over buildings against zero with a $t$-test and a sign test.
\item[H2] Under T1, the error of the forecast daily peak predicts the share better than the MAE (per-building Spearman correlations over forecaster-months, compared by a Wilcoxon signed-rank test).
\item[H3] Agreement between accuracy and value is higher under T2 than under T1 (paired Wilcoxon test over buildings).
\item[H4] Under T1 the reactive rule captures at least as large an annual share as the best forecaster in at least half of the buildings (sign test).
\end{itemize}
H2--H4 are adjusted with Holm's method. Secondary analyses (labelled as such) repeat H1 and H4 for the guarded controller (H1b, H4b) and compare the synthetic error types with Friedman tests. Robustness checks, all listed in the plan, use battery powers of 5\% and 20\% of the peak, demand charges of 10 and 30~\$/kW, RMSE in place of MAE, and the buildings without filled gaps.
